\chapter{Phân tích và thiết kế}
\label{chap:analysis-design}

\section{Yêu cầu chức năng}

Phần này nên liệt kê những chức năng hoặc nghiệp vụ chính mà sản phẩm cần đáp ứng.
Nếu hệ thống có nhiều vai trò, nên tách yêu cầu theo từng nhóm người dùng.

\section{Yêu cầu phi chức năng}

Những yêu cầu như hiệu năng, bảo mật, khả năng mở rộng, khả dụng và bảo trì thường nên đặt ở đây.
Nếu giảng viên yêu cầu tiêu chí cụ thể, nên đưa về dạng bảng để dễ đối chiếu.

\begin{table}[H]
  \centering
  \begin{tabularx}{0.92\linewidth}{|>{\bfseries}l|Y|}
    \hline
    Tiêu chí & Mô tả \\
    \hline
    Hiệu năng & Thời gian phản hồi mong muốn, số người dùng đồng thời, giới hạn tài nguyên. \\
    \hline
    Bảo mật & Xác thực, phân quyền, bảo vệ dữ liệu nhạy cảm, nhật ký hệ thống. \\
    \hline
    Vận hành & Cách triển khai, sao lưu, giám sát, và xử lý sự cố. \\
    \hline
  \end{tabularx}
  \caption{Mẫu bảng mô tả yêu cầu phi chức năng.}
  \label{tab:non-functional}
\end{table}

\section{Thiết kế giải pháp}

Từ bảng \ref{tab:non-functional}, nhóm có thể trình bày kiến trúc tổng thể, thành phần chính,
cơ sở dữ liệu, API, hay luồng xử lý. Nếu cần trích dẫn tài liệu kỹ thuật, có thể sử dụng
\verb|\cite{lamport1994latex}| như ví dụ \cite{lamport1994latex}.
